\section{Dirac-standard prose attack and heal}

\paragraph{Scope.}
This dossier attacks \texttt{proj.tex} at the level of sentence order,
phrasing, and epistemic status.  It uses
\texttt{agent\_material/09\_physics\_first\_rewrite\_patchset.tex} as the
physics-first input and writes no manuscript changes.  The standard is:
a sentence may stand only if it states the object, the operation, and the
status at the point where the derivation has earned them.

\paragraph{Dirac test.}
The sentence fails if it advertises a principle, says ``fixes'' without a
map, says ``same charge lattice'' before the charge projection is built,
uses ``physical'' for an automorphic or Gram-invariant object, or lets a
late comparison theorem act as an input.  The replacement sentence should
be shorter, should name the mathematical object, and should label
unproved physics as input, conjecture, or open problem.

\subsection{Sentence-level attack dossier}

\paragraph{D1. \texttt{proj.tex:56--57}.}
\textbf{Bad sentence.}
\begin{quote}\small
The operation studied here is the passage from a scalar protected
partition function to its automorphic square root.
\end{quote}
\textbf{Attack.}
This sentence begins with an operation whose physical input has not yet
been stated.  It also calls the scalar partition function protected before
the OP/DT/protected-index distinction is made.
\textbf{Replacement.}
\begin{quote}\small
For \(X=S\times E\), the proved scalar input is the
Oberdieck--Pixton primitive reduced curve-counting series, and the
automorphic square root is the normalized Borcherds determinant
\(\Delta_5\).
\end{quote}
\textbf{Status label.} Proved scalar input; constructed automorphic
determinant.

\paragraph{D2. \texttt{proj.tex:64--65}.}
\textbf{Bad sentence.}
\begin{quote}\small
The square-root object used here is not an analytic branch.  It is the
Borcherds determinant built from the virtual Ramond-sector K3 index.
\end{quote}
\textbf{Attack.}
The second sentence is close, but it suppresses the crucial half-index
status.  Dirac would name the virtual operation.
\textbf{Replacement.}
\begin{quote}\small
The square-root object is not an analytic branch; it is the Borcherds
determinant of the virtual class with superdimension
\(\frac12 Z_{\mathrm{K3}}=\phi_{0,1}\).
\end{quote}
\textbf{Status label.} Constructed virtual class; no Hilbert-space half.

\paragraph{D3. \texttt{proj.tex:70--71}.}
\textbf{Bad sentence.}
\begin{quote}\small
The Fock superdeterminant over the effective BPS charge cone is
\end{quote}
\textbf{Attack.}
The phrase ``BPS charge cone'' overstates the data.  At this point only
the product cone of Gram-invariant indices has been constructed.
\textbf{Replacement.}
\begin{quote}\small
The Fock superdeterminant over the cusp-positive Gram-index cone is
\end{quote}
\textbf{Status label.} Constructed product cone; physical charge cone
not proved.

\paragraph{D4. \texttt{proj.tex:81--87}.}
\textbf{Bad sentences.}
\begin{quote}\small
The duality principle fixes the rest of the paper.  The BPS charge
pairing has discriminant \(4nm-l^2\), and the degree map
\(\alpha(n,l,m)=2nf_2-lf_3+2mf_{-2}\) satisfies
\((\alpha(\gamma),\alpha(\gamma'))=-2\langle\gamma,\gamma'\rangle_{\mathrm{BPS}}\).
Hence the Weyl chamber and the Borcherds cone are the denominator
polarization of the same charge lattice, restricted to the active
Jacobi support, whose cusp polarization produced the determinant.
\end{quote}
\textbf{Attack.}
This is the central prose failure.  ``Principle fixes'' is rhetoric; the
pairing is called BPS before the physical charge lattice has been
constructed; ``same charge lattice'' hides the quotient from charges to
Gram invariants.
\textbf{Replacement.}
\begin{quote}\small
After the Harvey--Moore/DVV Gram-reduction input, the product variables
are indexed by \(\Gamma_{\mathrm{gram}}\cong\mathbb Z^3\) with
discriminant \(4nm-l^2\).  The map
\(\alpha(n,l,m)=2nf_2-lf_3+2mf_{-2}\) sends this index lattice to the
Lorentzian lattice used by the denominator identity.  The Weyl chamber is
therefore an automorphic chamber for the Gram-indexed product, not yet a
microscopic wall-crossing theorem.
\end{quote}
\textbf{Status label.} Assumption plus lattice calculation; physical
wall-crossing open.

\paragraph{D5. \texttt{proj.tex:89}.}
\textbf{Bad sentence.}
\begin{quote}\small
The automorphic square root also selects a denominator algebra.
\end{quote}
\textbf{Attack.}
``Selects'' is not a mathematical verb here.  The source is the
Gritsenko--Nikulin Borcherds product and denominator theorem.
\textbf{Replacement.}
\begin{quote}\small
The Gritsenko--Nikulin denominator theorem attaches to this product a
generalized Borcherds--Kac--Moody superalgebra.
\end{quote}
\textbf{Status label.} Cited denominator theorem.

\paragraph{D6. \texttt{proj.tex:110--114}.}
\textbf{Bad sentences.}
\begin{quote}\small
The associated current chiral algebra on the elliptic curve \(E\) is the
formal Beilinson--Drinfeld envelope of this denominator algebra.  The
remaining physical problem is to construct a microscopic BPS
factorization algebra for \(K3\times E\) and compare its protected
primitive Euler characteristics with this denominator algebra.
\end{quote}
\textbf{Attack.}
The first sentence says ``associated'' as if the current algebra is
geometric.  It is formal.  The second sentence should say what is missing:
states, bracket, and comparison.
\textbf{Replacement.}
\begin{quote}\small
One may formally form the Beilinson--Drinfeld current envelope on \(E\)
from the denominator algebra.  A physical BPS factorization algebra would
require a compactification construction of states, protected integration,
and a bracket whose primitive character and relations match
\(\mathfrak g_{\Delta_5}\).
\end{quote}
\textbf{Status label.} Formal construction; physical realization open.

\paragraph{D7. \texttt{proj.tex:123--125}.}
\textbf{Bad sentences.}
\begin{quote}\small
The square root is not first a product in a chosen coordinate system.
It is a duality-covariant automorphic section.  The product is its
trivialisation at the standard zero-dimensional cusp.
\end{quote}
\textbf{Attack.}
The sentences assert modular primacy before the physical charge/index
origin.  They are good after the Gram-reduction input, not before it.
\textbf{Replacement.}
\begin{quote}\small
Once the protected index has been reduced to Gram invariants, the
automorphic form is a section on the Siegel modular quotient.  Its product
expansion is the trivialization at the standard zero-dimensional cusp.
\end{quote}
\textbf{Status label.} Mathematical after physical input.

\paragraph{D8. \texttt{proj.tex:170--175}.}
\textbf{Bad sentences.}
\begin{quote}\small
The charge lattice is even, hence \((Q,Q)/2\) and \((P,P)/2\) are integral
and \((Q,P)\in\mathbb Z\).  The displayed matrix identity is immediate
from the definitions.  Its determinant is \(4nm-l^2\).  The last equality
is the definition of the Fourier character attached to the chemical
potential \(Z\).  This is the Harvey--Moore and
Dijkgraaf--Verlinde--Verlinde charge grading used by the BPS product
\(\cdots\).
\end{quote}
\textbf{Attack.}
The first and last sentences promote an assumption to a fact.  They also
conflate the compactification charge lattice with the product index.
\textbf{Replacement.}
\begin{quote}\small
By hypothesis the compactification charge lattice is even, so the three
Gram invariants \(n,l,m\) are integral.  The displayed matrix identity and
the determinant \(4nm-l^2\) are then immediate.  The Fourier monomial
\(q^nr^ls^m\) records this Gram matrix; the Harvey--Moore/DVV
interpretation of this grading is a physical input.
\end{quote}
\textbf{Status label.} Assumption plus calculation.

\paragraph{D9. \texttt{proj.tex:178--190}.}
\textbf{Bad sentence.}
\begin{quote}\small
The physical discriminant is
\(\Delta(\gamma)=\langle\gamma,\gamma\rangle_{\mathrm{BPS}}=4nm-l^2\).
\end{quote}
\textbf{Attack.}
The discriminant is a Gram discriminant at this point.  Calling it
physical presupposes the unproved charge-lattice theorem.
\textbf{Replacement.}
\begin{quote}\small
The Gram discriminant of the index triple is
\(\Delta_{\mathrm{gram}}(\gamma)=4nm-l^2\).
\end{quote}
\textbf{Status label.} Definition; physical interpretation assumed
elsewhere.

\paragraph{D10. \texttt{proj.tex:196}.}
\textbf{Bad sentence.}
\begin{quote}\small
It gives exactly the effective semigroup
\end{quote}
\textbf{Attack.}
``Exactly'' is emphasis without content.  The semigroup is determined by
the displayed cusp order.
\textbf{Replacement.}
\begin{quote}\small
This cusp order determines the semigroup
\end{quote}
\textbf{Status label.} Definition from order.

\paragraph{D11. \texttt{proj.tex:205--210}.}
\textbf{Bad sentences.}
\begin{quote}\small
The pairing \(\langle\ ,\ \rangle_{\mathrm{BPS}}\) is part of the charge
lattice.  The order on \(\Gamma_{\mathrm{eff}}\) is a Fock polarization at
one cusp.  A duality element transports this polarization to the
corresponding product chamber at the transported cusp; it does not
preserve the displayed inequalities inside the same coordinate chart.
\end{quote}
\textbf{Attack.}
The first sentence is false at the stated level of proof.  The pairing is
defined on the index lattice; its physical charge meaning is assumed.
\textbf{Replacement.}
\begin{quote}\small
The bilinear form is defined on the Gram-invariant index lattice.  Under
the Harvey--Moore/DVV input it is the image of the physical charge
pairing.  The order on \(\Gamma_{\mathrm{eff}}\) is a Fock polarization at
one cusp; a duality element transports the chamber, not the displayed
inequalities in the original coordinates.
\end{quote}
\textbf{Status label.} Definition plus physical input.

\paragraph{D12. \texttt{proj.tex:215--217}.}
\textbf{Bad sentence.}
\begin{quote}\small
The charge lattice is
\(\Gamma_{\mathrm{BPS}}=\mathbb Z^3,\ \gamma=(n,l,m),\ x^\gamma=q^nr^ls^m\).
\end{quote}
\textbf{Attack.}
This is the most direct overclaim.  \(\mathbb Z^3\) is the lattice of
Gram invariants used by the product, not the microscopic charge lattice.
\textbf{Replacement.}
\begin{quote}\small
The determinant is indexed by
\(\Gamma_{\mathrm{gram}}=\mathbb Z^3\), with
\(\gamma=(n,l,m)\) and \(x^\gamma=q^nr^ls^m\).
\end{quote}
\textbf{Status label.} Product indexing; microscopic charge lattice not
constructed.

\paragraph{D13. \texttt{proj.tex:400--403}.}
\textbf{Bad sentences.}
\begin{quote}\small
The determinant construction does not assume this object.  It uses only
the virtual classes \(\mathbb U_{n,l}\) and their signed dimensions.  The
algebraic object below is instead the Borcherds denominator algebra
selected by the same charge grading and the same signed multiplicities.
\end{quote}
\textbf{Attack.}
The last sentence reintroduces ``same charge grading'' after a paragraph
that correctly avoided Hilbert-space claims.
\textbf{Replacement.}
\begin{quote}\small
The determinant construction does not assume this object.  It uses only
the virtual classes \(\mathbb U_{n,l}\) and their signed dimensions.  The
algebraic object below is the Borcherds denominator algebra with the same
Gram-degree coefficients.
\end{quote}
\textbf{Status label.} Virtual determinant; denominator coefficients.

\paragraph{D14. \texttt{proj.tex:627--628}.}
\textbf{Bad sentence.}
\begin{quote}\small
Therefore the Oberdieck--Pixton monic product is exactly the square of
\end{quote}
\textbf{Attack.}
The word ``exactly'' is not needed.  The proof supplies the equality.
\textbf{Replacement.}
\begin{quote}\small
Therefore the Oberdieck--Pixton monic product is the square of
\end{quote}
\textbf{Status label.} Proved equality.

\paragraph{D15. \texttt{proj.tex:656--668}.}
\textbf{Bad sentence.}
\begin{quote}\small
For the D6--D2--D0 charge \(\gamma(h,d,n)\), the physical protected index
is expected to be represented by the Behrend-weighted Hilbert-scheme
invariant.
\end{quote}
\textbf{Attack.}
The status is correct but the grammar buries it.  Put the conjecture at
the front.
\textbf{Replacement.}
\begin{quote}\small
Conjecturally, the Behrend-weighted Hilbert-scheme invariant represents
the protected D6--D2--D0 index for the charge \(\gamma(h,d,n)\).
\end{quote}
\textbf{Status label.} Conjectural dictionary.

\paragraph{D16. \texttt{proj.tex:701--708}.}
\textbf{Bad sentence.}
\begin{quote}\small
The chosen square root is the Borcherds determinant already normalized as
\(\mathcal D_X=\Delta_5\), not an arbitrary analytic branch.
\end{quote}
\textbf{Attack.}
Good mathematics, but the sentence can be sharper by removing ``chosen''
and ``already''.
\textbf{Replacement.}
\begin{quote}\small
The square root used here is the normalized Borcherds determinant
\(\mathcal D_X=\Delta_5\), not an analytic branch.
\end{quote}
\textbf{Status label.} Definition and normalization.

\paragraph{D17. \texttt{proj.tex:840--842}.}
\textbf{Bad sentence.}
\begin{quote}\small
The geometric duality group is
\(G_{\mathrm{du}}=\PSp_4(\mathbb Z)\simeq\SO_+(\Lambda^{3,2})\).
\end{quote}
\textbf{Attack.}
Exterior-square linear algebra proves an arithmetic automorphic group.
It does not prove the physical duality group of the compactification.
\textbf{Replacement.}
\begin{quote}\small
The automorphic group used by the product is
\(\PSp_4(\mathbb Z)\simeq\SO_+(\Lambda^{3,2})\); its physical-duality
interpretation is part of the Harvey--Moore/DVV input.
\end{quote}
\textbf{Status label.} Proved arithmetic; physical interpretation
assumed.

\paragraph{D18. \texttt{proj.tex:866--870}.}
\textbf{Bad sentences.}
\begin{quote}\small
We now change polarization on the same BPS charge lattice.  The cusp
polarization produced the Fock determinant.  The denominator
polarization sends a charge \(\gamma=(n,l,m)\) to a Lorentzian degree
and makes the Borcherds chamber visible as the positive cone of that
grading.
\end{quote}
\textbf{Attack.}
The paragraph is too physical.  It assumes the same BPS lattice and calls
the index triple a charge.
\textbf{Replacement.}
\begin{quote}\small
We now pass from the Gram-index product cone to the Lorentzian lattice
used in the denominator identity.  The cusp polarization gives the Fock
determinant; the Lorentzian degree map gives the Borcherds chamber.
\end{quote}
\textbf{Status label.} Lattice construction; microscopic wall-crossing
not asserted.

\paragraph{D19. \texttt{proj.tex:1031--1038}.}
\textbf{Bad sentences.}
\begin{quote}\small
\texttt{\textbackslash begin\{definition\}[BPS charge degrees in the Weyl chamber]}
For \(\gamma=(n,l,m)\in\Gamma_{\mathrm{BPS}}\) set
\(\alpha(\gamma)=\alpha(n,l,m)=2nf_2-lf_3+2mf_{-2}\).
With \(z=z_1f_{-2}+z_2f_3+z_3f_2\), the character of the BPS charge is
\(x^\gamma=q^nr^ls^m=\exp(-\pi i(\alpha(\gamma),z))\).
\end{quote}
\textbf{Attack.}
The title and text import BPS language before the construction earns it.
\textbf{Replacement.}
\begin{quote}\small
\texttt{\textbackslash begin\{definition\}[Gram degrees in the Weyl chamber]}
For \(\gamma=(n,l,m)\in\Gamma_{\mathrm{gram}}\), set
\(\alpha(\gamma)=2nf_2-lf_3+2mf_{-2}\in\Lambda^{2,1}_{II}\).
With \(z=z_1f_{-2}+z_2f_3+z_3f_2\), the product monomial is
\(x^\gamma=q^nr^ls^m=\exp(-\pi i(\alpha(\gamma),z))\).
\end{quote}
\textbf{Status label.} Product character.

\paragraph{D20. \texttt{proj.tex:1046--1049}.}
\textbf{Bad sentences.}
\begin{quote}\small
Thus the Borcherds degree is the Lorentzian image of the BPS charge
grading.  The Weyl chamber below is determined by the cusp polarization
fixed in Section \ref{sec:duality-principle} together with the active
Jacobi support, since charges with zero exponent do not generate walls.
\end{quote}
\textbf{Attack.}
Again the phrase ``BPS charge grading'' overclaims.  The last clause
should say exponents, not charges.
\textbf{Replacement.}
\begin{quote}\small
Thus the Borcherds degree is the Lorentzian image of the Gram-index
grading.  The Weyl chamber below is determined by the cusp polarization
and the active Jacobi support; triples with zero exponent do not generate
walls of the denominator product.
\end{quote}
\textbf{Status label.} Denominator support.

\paragraph{D21. \texttt{proj.tex:1056--1065}.}
\textbf{Bad sentences.}
\begin{quote}\small
In the standard cusp chamber, the active charge support is
\(\Gamma_{\mathrm{act}}=\{\gamma=(n,l,m)\in\Gamma_{\mathrm{eff}}\mid
f(nm,l)\neq0\}\).  The inactive triples in
\(\Gamma_{\mathrm{eff}}\setminus\Gamma_{\mathrm{act}}\) are ordering data
only; their product exponents are zero.  The proof-grade Borcherds cone
is the real cone generated by the active positive BPS charges after the
Lorentzian degree map.
\end{quote}
\textbf{Attack.}
The computation is sound; the words ``charge'' and ``BPS charges'' are
not.
\textbf{Replacement.}
\begin{quote}\small
In the standard cusp chamber, the active Gram-index support is
\(\Gamma_{\mathrm{act}}=\{\gamma=(n,l,m)\in\Gamma_{\mathrm{eff}}\mid
f(nm,l)\neq0\}\).  The inactive triples have exponent zero.  The
Borcherds cone is the real cone generated by the active Gram degrees
\(\alpha(\Gamma_{\mathrm{act}})\).
\end{quote}
\textbf{Status label.} Computed support.

\paragraph{D22. \texttt{proj.tex:1211}.}
\textbf{Bad heading.}
\begin{quote}\small
\(\backslash\)section\{The denominator bracket\}
\end{quote}
\textbf{Attack.}
The section heading makes the bracket sound like the main object.  The
proved object is the denominator algebra; physical brackets are later.
\textbf{Replacement.}
\begin{quote}\small
\(\backslash\)section\{The denominator algebra\}
\end{quote}
\textbf{Status label.} Denominator-level theorem.

\paragraph{D23. \texttt{proj.tex:1218--1220}.}
\textbf{Bad sentence.}
\begin{quote}\small
This is the algebraic denominator square root.  It is the automorphic
correction of \(\bkm(A)\) constructed by Gritsenko--Nikulin,
with root spaces graded by the charge map.
\end{quote}
\textbf{Attack.}
``Charge map'' is too strong.  The grading map is the Gram-degree map.
\textbf{Replacement.}
\begin{quote}\small
This is the denominator algebra attached to the Borcherds product.  It is
the Gritsenko--Nikulin automorphic correction of \(\bkm(A)\),
with root spaces indexed by the Gram-degree map.
\end{quote}
\textbf{Status label.} Cited construction.

\paragraph{D24. \texttt{proj.tex:1231}.}
\textbf{Bad sentence.}
\begin{quote}\small
Harvey--Moore charge grading:
\end{quote}
\textbf{Attack.}
This label assigns physical meaning to the root-space superdimension.
\textbf{Replacement.}
\begin{quote}\small
Signed root-superdimension convention:
\end{quote}
\textbf{Status label.} Algebraic convention.

\paragraph{D25. \texttt{proj.tex:1236--1243}.}
\textbf{Bad sentences.}
\begin{quote}\small
For \(\gamma\in\Gamma_{\mathrm{BPS}}\) set
\(\mathcal A_{\mathrm{den},\gamma}:=\mathfrak g_{\Delta_5,\alpha(\gamma)}\).
If \(\alpha(\gamma)\) is not a root, this space is zero.  The bracket is
the Borcherds--Kac--Moody bracket of \(\mathfrak g_{\Delta_5}\).
\end{quote}
\textbf{Attack.}
The construction is algebraic and should not use \(\Gamma_{\mathrm{BPS}}\).
\textbf{Replacement.}
\begin{quote}\small
For \(\gamma\in\Gamma_{\mathrm{gram}}\), set
\(\mathcal A_{\mathrm{den},\gamma}:=\mathfrak g_{\Delta_5,\alpha(\gamma)}\).
If \(\alpha(\gamma)\) is not a root, this space is zero.  The bracket is
the Borcherds--Kac--Moody bracket of \(\mathfrak g_{\Delta_5}\).
\end{quote}
\textbf{Status label.} Algebraic root grading.

\paragraph{D26. \texttt{proj.tex:1254--1266}.}
\textbf{Bad sentences.}
\begin{quote}\small
\(\backslash\)label\{thm:denominator-algebra-identity\}
The denominator Lie superalgebra has the following properties.
\(\mathcal A_{\mathrm{den}}=\bigoplus_{\gamma\in\Gamma_{\mathrm{BPS}}}
\mathcal A_{\mathrm{den},\gamma}\).
Its signed charge multiplicities are the K3 half-index coefficients.
\end{quote}
\textbf{Attack.}
The label says bracket-level; the theorem is denominator-level.  The
grading is not \(\Gamma_{\mathrm{BPS}}\).
\textbf{Replacement.}
\begin{quote}\small
\(\backslash\)label\{thm:denominator-algebra-identity\}
The denominator Lie superalgebra has a
\(\Gamma_{\mathrm{gram}}\)-grading.  Its signed Gram-degree
multiplicities are the virtual K3 half-index coefficients.
\end{quote}
\textbf{Status label.} Denominator algebra theorem.

\paragraph{D27. \texttt{proj.tex:1323--1333}.}
\textbf{Bad sentences.}
\begin{quote}\small
\(\mathsf{Den}_{\Delta_5,E}\) is a formal chiral/factorization algebra on
\(E\) obtained from the denominator Lie superalgebra.  Its constant-mode
primitive Lie superalgebra is \(\mathcal A_{\mathrm{den}}\), and the
denominator of that primitive bracket is \(64^{-1}\Delta_5(2Z)\).  It is
not, without additional microscopic input, the BPS factorization algebra
of the compactification.
\end{quote}
\textbf{Attack.}
The first sentence is acceptable, but the theorem title ``Formal current
square root'' is too physical.  The third sentence should stand first in
the proof of status.
\textbf{Replacement.}
\begin{quote}\small
\(\mathsf{Den}_{\Delta_5,E}\) is the formal Beilinson--Drinfeld current
envelope of the denominator Lie superalgebra.  Its constant-mode
primitive Lie superalgebra is \(\mathcal A_{\mathrm{den}}\), whose
denominator is \(64^{-1}\Delta_5(2Z)\).  This formal current algebra is
not a BPS factorization algebra without a microscopic construction.
\end{quote}
\textbf{Status label.} Formal BD current envelope.

\paragraph{D28. \texttt{proj.tex:1343--1344}.}
\textbf{Bad sentence.}
\begin{quote}\small
The denominator statement is exactly Theorem
\ref{thm:denominator-algebra-identity}.
\end{quote}
\textbf{Attack.}
The word ``exactly'' is idle, and the referenced label is wrong in
status.
\textbf{Replacement.}
\begin{quote}\small
The denominator statement is Theorem
\ref{thm:denominator-algebra-identity}.
\end{quote}
\textbf{Status label.} Cross-reference hygiene.

\paragraph{D29. \texttt{proj.tex:1381--1384}.}
\textbf{Bad sentence.}
\begin{quote}\small
The protected charge lattice is
\(\Gamma_{\mathrm{BPS}}=\mathbb Z^3\), the completed positive support is
\(\Gamma_{\mathrm{eff}}\), and the Lorentzian degree is
\(\alpha(n,l,m)=2nf_2-lf_3+2mf_{-2}\).
\end{quote}
\textbf{Attack.}
This line is the Hall criterion's main overclaim.  The Hall object would
need a projection to the Gram lattice, not a declaration that the charge
lattice is \(\mathbb Z^3\).
\textbf{Replacement.}
\begin{quote}\small
The compactification charge lattice maps to
\(\Gamma_{\mathrm{gram}}\cong\mathbb Z^3\) with positive support
\(\Gamma_{\mathrm{eff}}\), and the projected degree is
\(\alpha(n,l,m)=2nf_2-lf_3+2mf_{-2}\).
\end{quote}
\textbf{Status label.} Required structure for a Hall realization.

\paragraph{D30. \texttt{proj.tex:1421--1424}.}
\textbf{Bad sentence.}
\begin{quote}\small
The first three conditions give a charge-graded Lie superalgebra whose
Grothendieck shadow has the same exponents as the Igusa Borcherds
product.  This is exactly the denominator/PBW shadow; by itself it is
only a character statement.
\end{quote}
\textbf{Attack.}
The first sentence should not say charge-graded.  The second sentence is
mostly right but should lose ``exactly''.
\textbf{Replacement.}
\begin{quote}\small
The first three conditions give a Gram-graded Lie superalgebra whose
Grothendieck shadow has the exponents of the Igusa Borcherds product.
This is the denominator/PBW character statement, not the bracket
identification.
\end{quote}
\textbf{Status label.} Character-level statement.

\paragraph{D31. \texttt{proj.tex:1446--1456}.}
\textbf{Bad sentences.}
\begin{quote}\small
The Oberdieck--Pixton scalar square fixes the protected character, and
the present paper constructs the Borcherds denominator bracket, but
neither result constructs the compact oriented critical Hall cosheaf, its
global protected integration map, the Hall structure constants, the
Borcherds relations inside the Hall super-commutator, the continuous Hopf
pairing, or the comparison between the compactification fields and the
Hall cosheaf.
\end{quote}
\textbf{Attack.}
OP fixes a scalar curve-counting series, not the protected character
without the physical dictionary.  ``Denominator bracket'' should be
``denominator algebra''.
\textbf{Replacement.}
\begin{quote}\small
Oberdieck--Pixton fix the primitive reduced scalar
\(-4096\Delta_5^{-2}\).  After the Harvey--Moore/DVV dictionary this is
expected to be the square of a protected character.  The present paper
constructs the Borcherds denominator algebra, but it does not construct
the compact oriented critical Hall cosheaf, protected integration, Hall
structure constants, Borcherds relations inside a Hall super-commutator,
continuous Hopf pairing, or comparison with compactification fields.
\end{quote}
\textbf{Status label.} OP theorem; physical character expected; Hall
realization open.

\paragraph{D32. \texttt{proj.tex:1459--1464}.}
\textbf{Bad sentences.}
\begin{quote}\small
The duality used here is automorphic BPS duality.  The same charge
lattice is read first at a cusp as a protected Fock determinant and then
in a Weyl chamber as a Borcherds denominator.  The formal current envelope
above is a chiral algebra, but it is not yet a physical factorization
algebra and not a chiral Koszul-duality theorem.
\end{quote}
\textbf{Attack.}
``Automorphic BPS duality'' and ``same charge lattice'' overstate.  The
good content is the last sentence.
\textbf{Replacement.}
\begin{quote}\small
The duality used here is automorphic duality after the protected-index
Gram reduction.  The same Gram-index lattice is read at the cusp as a
Fock determinant and in the Weyl chamber as a Borcherds denominator.  The
formal current envelope is a chiral algebra, but not a physical
factorization algebra or chiral Koszul-duality theorem.
\end{quote}
\textbf{Status label.} Automorphic theorem plus physical input.

\paragraph{D33. \texttt{proj.tex:1681--1693}.}
\textbf{Bad sentences.}
\begin{quote}\small
The equality \(\smult(\alpha(n,l,m))=f(nm,l)\) identifies the root-space
Euler characteristic of \(\mathcal A_{\mathrm{den}}\) with the
one-particle BPS index coefficient.  It does not prove an isomorphism
\(\mathfrak g_{\alpha(n,l,m)}\cong
\mathcal H^{\mathrm{BPS}}_{(n,l,m)}\) and it does not prove that the
microscopic operator product on BPS states is the Borcherds bracket.
\end{quote}
\textbf{Attack.}
The warning is correct but the first sentence still says BPS.  The
coefficient is the virtual half-index coefficient.
\textbf{Replacement.}
\begin{quote}\small
The equality \(\smult(\alpha(n,l,m))=f(nm,l)\) identifies the signed
root superdimension of \(\mathcal A_{\mathrm{den}}\) with the virtual K3
half-index coefficient.  It does not identify root spaces with
microscopic BPS state spaces, and it does not identify a microscopic
operator product with the Borcherds bracket.
\end{quote}
\textbf{Status label.} Signed dimension equality; state-space
identification open.

\paragraph{D34. \texttt{proj.tex:1711--1718}.}
\textbf{Bad sentences.}
\begin{quote}\small
If \(\mathcal B\) carries a bracket compatible with charge addition, the
determinant forgets that bracket and remembers exactly the class of
\(\mathcal B\) in \(K_0(\cdots)\).  Thus the determinant of
\(\mathcal A_{\mathrm{den}}=\mathfrak g_{\Delta_5}\) is the protected
shadow of its denominator bracket.  The bracket exists, but the
determinant sees only its charge-graded Grothendieck class.
\end{quote}
\textbf{Attack.}
The first sentence is good except for ``charge'' and ``exactly''.  The
second sentence says ``protected shadow'' without a protected object.
\textbf{Replacement.}
\begin{quote}\small
If \(\mathcal B\) carries a bracket compatible with the grading, the
determinant forgets that bracket and remembers the class of
\(\mathcal B\) in \(K_0\).  Thus the determinant of
\(\mathcal A_{\mathrm{den}}=\mathfrak g_{\Delta_5}\) is the
Gram-graded Euler shadow of its denominator bracket.
\end{quote}
\textbf{Status label.} Grothendieck-class statement.

\paragraph{D35. \texttt{proj.tex:1935--1943}.}
\textbf{Bad sentences.}
\begin{quote}\small
The congruence rows \(F_5,F_6,F_7,F_8\) have product-degree lattices in
Proposition~\ref{prop:congruence-product-degree-lattices}.  A nonabelian
Weyl--Borcherds promotion is additional structure: it is imported for
\(F_5,F_6,F_7,F_8\) from the CHL denominator theorem of Govindarajan.
\end{quote}
\textbf{Attack.}
This is mostly correct.  ``Promotion'' is rhetoric and hides the exact
data.
\textbf{Replacement.}
\begin{quote}\small
The congruence rows \(F_5,F_6,F_7,F_8\) have product-degree lattices in
Proposition~\ref{prop:congruence-product-degree-lattices}.  Their
nonabelian denominator data -- Weyl chamber, real roots, character, and
Weyl sum -- are imported from Govindarajan's CHL denominator theorem.
\end{quote}
\textbf{Status label.} Product data plus imported denominator theorem.

\paragraph{D36. \texttt{proj.tex:1948--1950}.}
\textbf{Bad sentences.}
\begin{quote}\small
For every row the product formula itself gives a charged primitive
product.  It gives a protected Fock determinant only after a one-particle
index and charge-lattice interpretation have been supplied.
\end{quote}
\textbf{Attack.}
``Charged'' and ``charge-lattice'' smuggle in physics.  The product gives
degrees.
\textbf{Replacement.}
\begin{quote}\small
For every row the product formula gives a primitive product indexed by
its product-degree lattice.  It gives a protected Fock determinant only
after a one-particle index and a physical charge projection have been
supplied.
\end{quote}
\textbf{Status label.} Product-level theorem; physical interpretation
conditional.

\paragraph{D37. \texttt{proj.tex:1968--1971}.}
\textbf{Bad sentence.}
\begin{quote}\small
If a future construction supplies the row by a CHL/K3 sector with
\(\phi^{(j)}=\tfrac12Z^{(j)}_{\mathrm{K3}}\), then the scalar-square
branch is \(Z_j=C_j/F_j^2\).
\end{quote}
\textbf{Attack.}
``Future construction'' is process language.  State the mathematical
condition.
\textbf{Replacement.}
\begin{quote}\small
If a CHL/K3 sector is constructed with
\(\phi^{(j)}=\tfrac12Z^{(j)}_{\mathrm{K3}}\), then the corresponding
scalar-square branch has the form \(Z_j=C_j/F_j^2\).
\end{quote}
\textbf{Status label.} Conditional statement.

\paragraph{D38. \texttt{proj.tex:2003--2008}.}
\textbf{Bad sentences.}
\begin{quote}\small
The physical square-root atlas used in Vol.~III has weights
\((5,2,1,1,\tfrac12,1,\tfrac14,0)\).  That line is not the table of the
theorem above unless a row map from the Clery--Gritsenko forms to the
physical rows has first been built.
\end{quote}
\textbf{Attack.}
``Physical square-root atlas'' is not a proved object in this paper.
The row-map obstruction is correct and should lead.
\textbf{Replacement.}
\begin{quote}\small
Vol.~III uses a separate comparison notation with weights
\((5,2,1,1,\tfrac12,1,\tfrac14,0)\).  It is not the
Clery--Gritsenko table unless a row map to the forms \(F_j\) has been
constructed.
\end{quote}
\textbf{Status label.} Comparison notation; row map open.

\paragraph{D39. \texttt{proj.tex:2380--2382}.}
\textbf{Bad sentence.}
\begin{quote}\small
The displayed sets are exactly
\(P_{t,II,\overline0}(\mathcal M_{t,II,\overline0})\) for \(t=2,3,4\).
\end{quote}
\textbf{Attack.}
Here ``exactly'' is acceptable only if it is doing citation work.  Dirac
would remove it.
\textbf{Replacement.}
\begin{quote}\small
The displayed sets are
\(P_{t,II,\overline0}(\mathcal M_{t,II,\overline0})\) for \(t=2,3,4\).
\end{quote}
\textbf{Status label.} Cited identification.

\paragraph{D40. \texttt{proj.tex:2461--2462}.}
\textbf{Bad sentence.}
\begin{quote}\small
Their formula (5.1) is exactly the displayed Weyl--Kac--Borcherds
identity.
\end{quote}
\textbf{Attack.}
Again, ``exactly'' is idle.
\textbf{Replacement.}
\begin{quote}\small
Their formula (5.1) is the displayed Weyl--Kac--Borcherds identity.
\end{quote}
\textbf{Status label.} Source identification.

\paragraph{D41. \texttt{proj.tex:2822--2830}.}
\textbf{Bad sentences.}
\begin{quote}\small
In particular the product already fixes the signed root multiplicity
function \(\mu_j\) and the product-prefactor degree \(\rho_j^{\mathrm{prod}}\),
but it does not fix \(P_j,\mathcal M_j,W_j,\epsilon_j\), or the imaginary
simple-root coefficients \(N_j(a)\).
\end{quote}
\textbf{Attack.}
Good claim, but ``signed root multiplicity'' should be ``product
exponent'' until the denominator data are present.
\textbf{Replacement.}
\begin{quote}\small
The product fixes the exponent function \(\mu_j\) and the prefactor
degree \(\rho_j^{\mathrm{prod}}\).  It does not fix the real roots, Weyl
chamber, Weyl group, reflection character, or imaginary simple-root
coefficients.
\end{quote}
\textbf{Status label.} Product-level data only.

\paragraph{D42. \texttt{proj.tex:3235--3240}.}
\textbf{Bad sentences.}
\begin{quote}\small
Thus all eight rows now have formal product/current-envelope objects.
Rows \(1\)--\(4\) have the Gritsenko--Nikulin denominator algebras.
Rows \(5\)--\(8\) have the Govindarajan CHL/WKB denominator algebras in
the imported normalization.  The row \(F_7\) is imported only on the
\(\mu_4\)-multiplier cover; it supplies no Vol.~III physical boundary
theorem.
\end{quote}
\textbf{Attack.}
``Now'' is process prose.  ``Current-envelope objects'' should not be
automatic from products unless the formal construction has been stated
row by row.
\textbf{Replacement.}
\begin{quote}\small
Each row has the product-level formal object specified by its local data.
Rows \(1\)--\(4\) have Gritsenko--Nikulin denominator algebras.  Rows
\(5\)--\(8\) have Govindarajan CHL/WKB denominator algebras in the
imported normalization.  The row \(F_7\) is imported only on the
\(\mu_4\)-multiplier cover and supplies no Vol.~III physical boundary
theorem.
\end{quote}
\textbf{Status label.} Local formal/product status; imported denominator
status.

\paragraph{D43. \texttt{proj.tex:3267--3276}.}
\textbf{Bad sentences.}
\begin{quote}\small
The Vol.~III physical square-root notation
\(\Phi^{\mathrm{V3}}_1,\ldots,\Phi^{\mathrm{V3}}_8\), with weights
\((5,2,1,1,\tfrac12,1,\tfrac14,0)\), is a different row system until a
row map has been constructed.  In that physical row system, the integral
rows are \(\cdots\) and the boundary rows are \(\cdots\).
\end{quote}
\textbf{Attack.}
``Physical row system'' overstates.  The section only compares notation.
\textbf{Replacement.}
\begin{quote}\small
The Vol.~III comparison notation
\(\Phi^{\mathrm{V3}}_1,\ldots,\Phi^{\mathrm{V3}}_8\), with weights
\((5,2,1,1,\tfrac12,1,\tfrac14,0)\), is a different row system until a
row map has been constructed.  In that comparison system, the integral
and boundary labels are as displayed below.
\end{quote}
\textbf{Status label.} Notational comparison; row map open.

\paragraph{D44. \texttt{proj.tex:3283}.}
\textbf{Bad sentence.}
\begin{quote}\small
For the physical rows the surviving boundary problem is
\end{quote}
\textbf{Attack.}
``Surviving'' is audit prose.  ``Physical rows'' should be ``comparison
rows'' unless the physical host is constructed.
\textbf{Replacement.}
\begin{quote}\small
For the Vol.~III comparison rows the open boundary cases are
\end{quote}
\textbf{Status label.} Open comparison problem.

\paragraph{D45. \texttt{proj.tex:3401--3406}.}
\textbf{Bad sentences.}
\begin{quote}\small
The following compatibility dictionary is restricted to the \(N=1\) Igusa
boundary.  The dictionary itself supplies no row map
\(\Phi^{\mathrm{V3}}_j\longleftrightarrow F_i\)
and no Vol.~III physical row-map denominator theorem for any congruence
row.
\end{quote}
\textbf{Attack.}
This is mostly correct, but ``compatibility dictionary'' is too broad for
material that is only a local comparison.
\textbf{Replacement.}
\begin{quote}\small
The following comparison is restricted to the \(N=1\) Igusa boundary.  It
supplies no row map
\(\Phi^{\mathrm{V3}}_j\longleftrightarrow F_i\)
and no physical row-map denominator theorem for any congruence row.
\end{quote}
\textbf{Status label.} Restricted comparison.

\subsection{Theorem-order changes}

\begin{enumerate}
\item Put the OP scalar theorem before the duality section.  The reader
should see the proved scalar identity before seeing \(\Sp_4\),
\(\Lambda^{3,2}\), or a Weyl chamber.
\item Put the K3 Ramond index and virtual half-index immediately after
the scalar theorem.  This proves that the square root is virtual before
the denominator algebra is introduced.
\item Put the Gram-reduction assumption after the virtual half-index.
This is the first legitimate use of the charge pair \((Q,P)\).
\item Put exterior-square arithmetic after the Gram-reduction assumption.
Then \(\PSp_4(\mathbb Z)\simeq\SO_+(\Lambda^{3,2})\) is an automorphic
group, with physical-duality status explicitly assumed.
\item Put the Weyl chamber after the exterior-square calculation.  Its
positive cone is an automorphic denominator chamber, not the origin of
the charge lattice.
\item Put the denominator algebra after the Weyl chamber and before any
Hall language.  The denominator identity supplies signed
superdimensions, not microscopic states.
\item Put the Hall/factorization criterion after the denominator algebra.
This is where a physical BPS bracket may be named, with open-problem
status.
\item Keep the Clery--Gritsenko and Vol.~III rows late.  They are
comparison and extension material, not input to the \(N=1\) derivation.
\end{enumerate}

\subsection{Claims strengthened}

\begin{enumerate}
\item The OP identity becomes a proved scalar theorem, not a protected
character theorem.
\item \(\Delta_5\) becomes the normalized Borcherds determinant of a
virtual half-index, not an analytic branch or Hilbert-space square root.
\item \(\Gamma_{\mathrm{gram}}\) replaces \(\Gamma_{\mathrm{BPS}}\) until
the compactification charge projection is constructed.
\item The \(\PSp_4\simeq\SO(3,2)\) statement is kept as arithmetic; its
physical-duality role is marked as input.
\item The Weyl chamber is a denominator chamber; microscopic
wall-crossing remains open.
\item The BKM object is a denominator algebra; the physical BPS bracket
requires Hall/factorization data.
\item Late row material is quarantined as comparison or imported
denominator material.
\end{enumerate}

\subsection{Claims refused}

\begin{enumerate}
\item Refused: ``the duality principle fixes the rest of the paper.''
\item Refused: ``\(\mathbb Z^3\) is the BPS charge lattice.''
\item Refused: ``the physical discriminant'' before the charge projection.
\item Refused: ``same BPS charge lattice'' in the Weyl transition.
\item Refused: ``OP fixes the protected character.''
\item Refused: ``denominator bracket'' as the name for the proved object.
\item Refused: ``product fixes signed root multiplicities'' before
denominator data.
\item Refused: Vol.~III comparison notation as a physical row system.
\end{enumerate}

\subsection{Commands run}

\begin{verbatim}
sed -n '1,220p' AGENTS.md
sed -n '1,220p' CLAUDE.md
sed -n '1,220p' /Users/raeez/ecosystem/INVARIANTS.md
sed -n '1,220p' /Users/raeez/ecosystem/AGENTS-HARNESS.md
test -e agent_material/10_dirac_prose_attack_heal.tex && \
  sed -n '1,240p' agent_material/10_dirac_prose_attack_heal.tex || true
sed -n '1,240p' agent_material/09_physics_first_rewrite_patchset.tex
sed -n '240,520p' agent_material/09_physics_first_rewrite_patchset.tex
sed -n '520,700p' agent_material/09_physics_first_rewrite_patchset.tex
wc -l proj.tex agent_material/09_physics_first_rewrite_patchset.tex proj.bib
rg -n -e "^This |^Thus |^We |^The point|not merely|exactly|canonical|..." \
  proj.tex
nl -ba proj.tex | sed -n '55,225p'
nl -ba proj.tex | sed -n '250,410p'
nl -ba proj.tex | sed -n '620,675p'
nl -ba proj.tex | sed -n '688,736p'
nl -ba proj.tex | sed -n '730,875p'
nl -ba proj.tex | sed -n '1028,1088p'
nl -ba proj.tex | sed -n '1128,1150p'
nl -ba proj.tex | sed -n '1210,1475p'
nl -ba proj.tex | sed -n '1678,1722p'
nl -ba proj.tex | sed -n '1890,2015p'
nl -ba proj.tex | sed -n '2368,2490p'
nl -ba proj.tex | sed -n '2816,2830p'
nl -ba proj.tex | sed -n '2898,2965p'
nl -ba proj.tex | sed -n '3060,3220p'
nl -ba proj.tex | sed -n '3220,3410p'
python3 compute/verify_lattice.py
python3 compute/verify_square_root.py
rg -n -F -e "The duality principle fixes" \
  -e "effective BPS charge cone" \
  -e "same BPS charge lattice" \
  -e "The geometric duality group is" \
  -e "protected character" \
  -e "The charge lattice is" \
  -e "The physical discriminant" \
  -e "BPS charge degrees" \
  -e "positive BPS charges" \
  -e "The protected charge lattice is" \
  -e "one-particle BPS" \
  -e "physical square-root atlas" \
  -e "physical row system" proj.tex
git status --short
perl -ne 'print "$.:$_" if /[^\x00-\x7F]/' \
  agent_material/10_dirac_prose_attack_heal.tex
chktex -q agent_material/10_dirac_prose_attack_heal.tex || true
sed -n '350,460p' agent_material/10_dirac_prose_attack_heal.tex
sed -n '900,1010p' agent_material/10_dirac_prose_attack_heal.tex
git status --short -- agent_material/10_dirac_prose_attack_heal.tex
\end{verbatim}

\subsection{Residual prose obligations}

\begin{enumerate}
\item After applying this dossier, run a zero-hit scan for
\(\Gamma_{\mathrm{BPS}}=\mathbb Z^3\), ``same BPS charge lattice'',
``protected character'' near the OP theorem, ``geometric duality group'',
and ``physical row system''.
\item Recheck every occurrence of ``exactly''.  Keep it only when it
marks equality with a cited formula or a displayed computation.
\item Recheck every occurrence of ``physical''.  It must refer either to
an explicit physical assumption, a conjectural dictionary, or an open
construction.
\item Recheck every theorem title containing ``square root''.  The title
must say whether the square root is scalar, virtual determinant,
denominator, or microscopic.
\end{enumerate}
